%======================================================================
%  Companion to Chapter 4 -- Plasticity: yield, flow and GSM
%======================================================================
\chapter{Plasticity: yield, flow and generalized standard materials}\label{cc:ch4}

\framepart{Outline}

Chapter~\ref{B-ch:pl} of the notes solves three problems in closed form: the
filament, the sheet stretched in plane strain, the thick-walled cylinder.
\op{pasapas}, the incremental procedure of Cast3M, solves them with its own
return maps. This chapter uses it as a check of the closed forms, and the
closed forms as a check of what Cast3M calls $H$, \texttt{TRAC} or
\texttt{EPSE}. You should then be able to set up a \op{pasapas} computation
with an imposed displacement or pressure history, extract a stress history at
a Gauss point, and read a residual stress field.

\section{Where Cast3M enters Chapter 4}\label{cc:sec4-map}
\index{Cast3M!Chapter 4}%

\begin{gbox}{Chapter 4 of the notes and Cast3M}\label{cc:tab4}
\small
\begin{tabular}{@{}p{42mm}@{\hspace{3mm}}p{98mm}@{}}
Filament, Sec.~\ref{B-sec:pl-1d}; Ex.~\ref{B-exo:pl-cyclic}
 & \textbf{Exercise \ref{exo:cc-filament}}: one bar element, isotropic (\texttt{TRAC}) and
   kinematic (\texttt{H}) hardening, turning-point stresses.\\[1mm]
Sheet, Ex.~\ref{B-exo:pl-sheet}
 & \textbf{Exercise \ref{exo:cc-sheet}}: one element in plane strain, first yield
   $\sY/\sqrt{1-\nu+\nu^2}$, path to $(2\sY/\sqrt3,\ \sY/\sqrt3)$.\\[1mm]
Cylinder, Ex.~\ref{B-exo:pl-cyl}
 & \textbf{Exercise \ref{exo:cc-residual}}: loading to a front $c=1.5a$, unloading, residual
   hoop stresses $-0.665k$ and $0.146k$.\\[1mm]
Tension--torsion, Ex.~\ref{B-exo:pl-tt}
 & To do: a thin tube (axisymmetric with torsion, Fourier mode~0), radial and
   non-radial paths in $(\sigma,\tau)$, first yield on $\sigma^2+3\tau^2=\sY^2$.\\[1mm]
Torsion, bending, Ex.~\ref{B-exo:pl-torsion}, \ref{B-exo:pl-bending}
 & To do: limit ratios $4/3$ and $3/2$, springback and residual stresses
   ($-7k/24$ at the surface from $c=R/2$; $\mp\sY/2$).\\[1mm]
Viscoplastic filament, Ex.~\ref{B-exo:pl-perzyna}
 & With the hand-written return map of Chapter~\ref{cc:ch5} (entry
   \texttt{ETA}): relaxation to $251.01$\,MPa at $5\tau$.\\[1mm]
Drucker--Prager, Ex.~\ref{B-exo:pl-dp}
 & To do: the \texttt{DRUCKER\_PRAGER} model of \op{pasapas}, a triaxial
   test, the apex.\\
\end{tabular}
\end{gbox}

\section{Reading a \texttt{pasapas} computation}\label{cc:sec4-pasapas}
\index{pasapas@\texttt{pasapas}}%
\begin{castemcom}
\op{pasapas} reads a table: the model, the material, the constraints, the
loading (a spatial field times a time history, built by \op{char}), the list
of instants. It writes its results in the same table: \texttt{DEPLACEMENTS},
\texttt{CONTRAINTES}, \texttt{VARIABLES\_INTERNES}, indexed by the number of
the instant from 0.

A Gauss-point value is extracted by zone, element and point; the cumulated
plastic strain is the internal variable \texttt{EPSE}, and its mask
\texttt{masq ep 'SUPERIEUR' 0.} marks the plastic zone.
\end{castemcom}
\begin{castemlines}
\begin{verbatim}
tab = table;
tab . 'MODELE' = mo;
tab . 'CARACTERISTIQUES' = ma;
tab . 'BLOCAGES_MECANIQUES' = aa;
tab . 'CHARGEMENT' = char 'MECA' fp ev;
tab . 'TEMPS_CALCULES' =
      prog 0. pas 0.05 2.;
pasapas tab;

sg = tab . 'CONTRAINTES' . 40;
s1 = extr sg 'SMTT' 1 1 1;
ep = exco
  (tab . 'VARIABLES_INTERNES' . 20)
  'EPSE' 'SCAL';
\end{verbatim}
\end{castemlines}

\framepart{Summary}

\begin{gbox*}
\begin{itemize}
\item \textbf{Homogeneous tests need one element}: the sheet and the
  filament do not depend on the mesh, only on the steps, and the return map is
  exact on monotone proportional steps.
\item \textbf{The closed forms calibrate the code}: the reverse yield stress
  of the kinematic filament tells which modulus Cast3M calls $H$.
\item \textbf{Residual stresses are a difference}: load, unload elastically,
  read the stresses at the last instant.
\end{itemize}
\end{gbox*}

\framepart{Exercises}

\exercise{The sheet in plane strain with \texttt{pasapas}}\label{exo:cc-sheet}
\index{sheet in plane strain!Cast3M}%
Compute the sheet of Exercise~\ref{B-exo:pl-sheet} ($E=200$\,GPa, $\nu=0.3$,
$\sY=250$\,MPa, perfect plasticity) with one element, stretched along $x$ to
$\varepsilon_{xx}=0.004$ in 400 steps, the faces $y=\pm h$ free. Compare the
first yield and $\sigma_{xx}$, $\sigma_{zz}$ at four strains with the closed
form, and plot the loading path in $(\sigma_{xx},\sigma_{zz})$. The course file
\texttt{sheet\_plane\_strain.dgibi}, in SI units with $\sY=200$\,MPa and the
stretching along $y$, does the same computation.
\begin{solution}
First yield at $\sigma_{xx}=\sY/\sqrt{1-\nu+\nu^2}=281.272$\,MPa. At
$\varepsilon_{xx}=0.002$, 0.0025, 0.003, 0.004 the closed form
(\texttt{exo\_sheet.py}) gives $\sigma_{xx}=286.484$, 287.765, 288.303,
288.615 and $\sigma_{zz}=112.499$, 124.049, 131.470, 139.197\,MPa, on the way
to the limits $2\sY/\sqrt3=288.675$ and $\sY/\sqrt3=144.338$. The path in
$(\sigma_{xx},\sigma_{zz})$ climbs along the von Mises ellipse
$\sigma_{xx}^2-\sigma_{xx}\sigma_{zz}+\sigma_{zz}^2=\sY^2$ towards the point
where the normal is parallel to $\vect{e}_x$. With the data of the course file
($\sY=200$\,MPa): first yield 225.02\,MPa, limits 230.94 and 115.47\,MPa.
\untested{\texttt{cast3m/ch4/sheet\_vm\_250.dgibi}}
\lstinputlisting[style=gibstyle,firstline=23]{cast3m/ch4/sheet_vm_250.dgibi}
\end{solution}

\exercise{The filament under a strain cycle}\label{exo:cc-filament}
\index{Bauschinger effect!Cast3M}%
One bar element, $E=200$\,GPa, $\sY=250$\,MPa, follows the strain path
$0\to3\varepsilon_Y\to-3\varepsilon_Y\to3\varepsilon_Y$. Compare the stresses at
the turning points with Exercise~\ref{B-exo:pl-cyclic} for (i) isotropic
hardening $K=3$\,GPa, given to Cast3M as a tensile curve, and (ii) kinematic
hardening $H=3$\,GPa.
\begin{solution}
(i)~$257.389$, $-271.949$, $286.079$\,MPa; (ii)~$257.389$, $-257.389$,
$257.389$\,MPa, independent of the number of steps. The tensile curve of (i)
has the elastoplastic slope $EK/(E+K)$ after $\varepsilon_Y$
(Section~\ref{B-ssec:pl-1d-iso}). Case~(ii) is a calibration, decided by
the first value: $257.389$\,MPa if Cast3M's $H$ is the uniaxial plastic
modulus of the notes ($q=H\varepsilon^p$), $254.95$ if that modulus is
$\tfrac23H$, $261.00$ if it is $\tfrac32H$. The later turning points are
symmetric whatever the modulus (the reverse yield itself occurs at
$257.389-2\sY=-242.61$\,MPa). If the convention differs, $H$ must be rescaled
in every kinematic computation (Exercises~\ref{exo:cc-bree} and \ref{exo:cc-id-sweep}--\ref{exo:cc-id-jax}).
\untested{\texttt{cast3m/ch4/filament\_cyclic.dgibi}}
\lstinputlisting[style=gibstyle,firstline=26]{cast3m/ch4/filament_cyclic.dgibi}
\end{solution}

\exercise{Residual stresses in the cylinder}\label{exo:cc-residual}
\index{thick-walled cylinder!residual stresses}%
The cylinder of Exercise~\ref{B-exo:pl-cyl} ($a=100$, $b=200$\,mm, plane
strain, perfect plasticity, $\sY=250$\,MPa) is loaded to the pressure that
puts the plastic front at $c=1.5a$, then unloaded. Compute the front at the
peak and the residual stresses, and compare with the closed form.
\begin{solution}
With $k=2\sY/\sqrt3=288.68$\,MPa: $p_Y=108.25$\,MPa, the front reaches
$c=1.5a$ at $p=180.20$\,MPa, $p_L=200.09$\,MPa. The unloading is elastic
($180.20<2p_Y$). Residual hoop stress $-0.6646k=-191.85$\,MPa at the bore and
$+0.1464k=42.25$\,MPa at the outer surface, radial stress zero on both faces.
Expect the computed front slightly beyond $1.5a$ and residual stresses a few
MPa off: the closed form takes $\sigma_\theta-\sigma_r=k$, exact for von Mises
only with incompressible plastic flow (solution of
Exercise~\ref{B-exo:pn-fem}).
\untested{\texttt{cast3m/ch4/cylinder\_residual.dgibi}}
\lstinputlisting[style=gibstyle,firstline=24]{cast3m/ch4/cylinder_residual.dgibi}
\end{solution}

\begin{chapterbib}{9}
\bibitem{cc-pasapas} Cast3M, notice of the procedure \texttt{PASAPAS} and
  of the operator \texttt{MATE}, \url{https://www-cast3m.cea.fr}.
\end{chapterbib}
